\chapter*{How to Use This Book}
\addcontentsline{toc}{chapter}{How to Use This Book}

Each chapter follows the same pattern:

\begin{enumerate}
  \item a motivating question;
  \item definitions and theorems;
  \item small hand computations;
  \item a short \Python/\NumPy experiment;
  \item a summary and forward links;
  \item exercises divided into Theory, By Hand, Python, and Challenge.
\end{enumerate}

The exercises are divided by purpose.

\begin{description}
  \item[Theory.] Proofs, definitions, counterexamples, and conceptual transfer.
  \item[By Hand.] Small mechanical computations. These build intuition for
    algorithms without turning the work into arithmetic endurance.
  \item[Python / NumPy.] Short programs using basic \Python\ and \NumPy. These
    show how the same ideas scale on a computer and how library routines relate
    to algorithms.
  \item[Challenge.] Optional or cumulative problems, clearly marked.
\end{description}

Hints for selected exercises appear near the end of the book. A hint gives the
first useful move, not a full solution. Problems with hints are still meant to
be solved by the student.

The computational labs are larger than the Python exercises. A Python exercise
usually tests one idea from the chapter. A lab is a guided investigation using
several ideas, often with plots, images, or richer numerical experiments.
